\section{Projective realization and removal of the surface blowups}
\label{sec:projective}

Theorem~\ref{thm:orbifold-cyclic} supplies a map $Y\to\cS$, where $Y$
is obtained by point blowups from the simple abelian surface $A$, and
where the image of $\pi_1(Y)$ is infinite cyclic.  We now make two changes
of ambient space.  The first replaces the orbifold by a smooth projective
variety containing $Y$.  The second removes the point blowups from $Y$
while preserving the homomorphism on fundamental groups.  Together they
produce the embedding required by Proposition~\ref{prop:large-reduction}.

\subsection{Replacing the orbifold ambient space}

The projective-bundle construction below is a relative version of the
free-locus method in the Godeaux--Serre construction; compare
\cite[\S20]{Serre1958} and \cite[Remark~1.4 and \S5]{Totaro1999}.
We include the containment and fundamental-group arguments, since the
specified surface must survive the replacement.

\begin{lemma}\label{lem:ambient-realization}
Let $\mathcal T$ be a smooth connected effective proper complex
Deligne--Mumford stack with projective coarse space.  Let $Y$ be a smooth
projective surface and let $f:Y\to\mathcal T$ have image in the locus
with trivial stabilizers.  There exist a smooth connected projective
variety $S$, an embedding $i:Y\hookrightarrow S$, and a morphism
$h:S\to\mathcal T$ such that $h\circ i=f$ and
\[
 h_*:\pi_1(S)\xrightarrow{\ \sim\ }\pi_1(\mathcal T).
\]
One may take $\dim S=5$.
\end{lemma}

\begin{proof}
We first pass to a projective bundle whose nonordinary locus has
arbitrarily large codimension, without changing the fundamental group.
Choose an embedding $Y\hookrightarrow\PP^{b-1}$ and form
\[
 \mathcal E=(\mathcal O_{\mathcal T}\oplus T_{\mathcal T})^{\oplus a}
                 \oplus\mathcal O_{\mathcal T}^{\oplus b},
 \qquad \mathcal P=\PP_{\mathcal T}(\mathcal E),
\]
using the convention that projectivization parametrizes lines.  The
trivial summands give a closed substack
$\mathcal T\times\PP^{b-1}\subset\mathcal P$, in which the graph of
$f$ is a closed copy of $Y$.  This copy lies in the ordinary locus,
because $f(Y)$ does.

The representation of each stabilizer on the tangent space of
$\mathcal T$ is faithful; in particular, every nonidentity element acts
nontrivially.  Indeed a finite holomorphic action is linearizable at a
fixed point, and effectiveness rules out an element with trivial tangent
action.  The added trivial line ensures
that its action on $\mathcal O\oplus T_{\mathcal T}$ is not scalar.
Consequently every eigenspace in $\mathcal E$ has codimension at least
$a$, including after the extra trivial summands are added.  The
projective fixed locus of each such element is a union of projective
eigenspaces.  It follows, in local finite quotient charts, that the
nonordinary locus of $\mathcal P$ has codimension at least $a$.
The same bound holds in its coarse space $P$.

The coarse space $P$ is projective.  Stabilizer orders are bounded on
the proper stack, so a sufficiently divisible power of the tautological
polarization has trivial stabilizer actions and
descends to $P$; it is relatively ample over the coarse space of
$\mathcal T$.  These statements can be checked on finite quotient
charts, where descent is precisely triviality of the stabilizer action
on the fibers of the line bundle.  Twisting by a sufficiently ample
bundle from the projective base gives a polarization of $P$.
Let $U\subset P$ be the ordinary locus.  It is smooth and identifies
with the ordinary locus of $\mathcal P$.

Choose $a>5$.  Removing a subset of complex codimension at least two
from a smooth orbifold does not change its orbifold fundamental group:
paths and their homotopies can be moved off that subset in smooth
charts, or in a smooth frame presentation.  Also, the projective-space
bundle $\mathcal P\to\mathcal T$ induces an isomorphism on fundamental
groups, since its fibers are simply connected.  Thus the natural maps
give
\begin{equation}\label{proj:ordinary-group}
 \pi_1(U)\simeq\pi_1(\mathcal P)\simeq\pi_1(\mathcal T).
\end{equation}

It remains to cut down inside $U$ while retaining the embedded surface.
Put $N=\dim P$.  In a sufficiently high projective embedding of $P$,
take $N-5$ general hyperplane sections through $Y$.  Their linear systems
are free off $Y$ and generate the normal differentials along $Y$.
The locus where $r$ general normal differentials fail to be independent
has expected codimension
\begin{equation}\label{proj:rank-loss}
 (N-2)-r+1
\end{equation}
on the surface.  For $r=N-5$ this is $4>2$, so the final intersection
is smooth along $Y$.  Bertini gives smoothness elsewhere in $U$.
Since $P\setminus U$ has codimension greater than five and is disjoint
from $Y$, a general constrained intersection avoids it.  We obtain a
smooth projective fivefold $S\subset U$ containing $Y$.

To verify the assertion about $\pi_1$, first take unconstrained general
hyperplane sections in this same embedding.  The quasi-projective
Lefschetz theorem of Hamm and L\^e
\cite[Theorem~1.1.3(ii)]{HammLe} says that a smooth open variety of
complex dimension $n$ is obtained, up to homotopy, from a general
hyperplane section by attaching cells of dimensions at least $n$.
At every stage here the dimension before cutting is at least six.
The successive inclusions therefore preserve connectedness and induce
isomorphisms on fundamental groups.  The final intersection avoids
$P\setminus U$, by the same dimension count.

Now consider the full product of the hyperplane parameter spaces,
without a containment requirement.  The tuples whose intersections
have dimension five, are smooth, and avoid $P\setminus U$ form a
nonempty Zariski-open set $\Omega$.  It contains both a general
unconstrained tuple and the constrained tuple just constructed, and it
is path connected.  The universal complete intersection over $\Omega$
is a smooth proper family with a map to $U$.  Along a path in $\Omega$,
Ehresmann's theorem identifies its fibers by diffeomorphisms and gives
a homotopy of their inclusions into $U$.  Thus the isomorphism on
fundamental groups for an unconstrained fiber also holds for $S$.
Composing $S\hookrightarrow U\subset\mathcal P$ with the bundle
projection gives $h$, with $h\circ i=f$, and
\eqref{proj:ordinary-group} proves the lemma.
\end{proof}

Applied to $Y\to\cS$, this gives a smooth connected projective ambient
variety $S$ with an embedding of $Y$ and
\begin{equation}\label{proj:blownup-image}
 \im\bigl(\pi_1(Y)\longrightarrow\pi_1(S)\bigr)\simeq\Z.
\end{equation}
The remaining issue is that Proposition~\ref{prop:large-reduction}
requires $A$ itself.  We address one point blowup at a time.

\subsection{Realizing a surface blowdown in a new ambient variety}

\begin{lemma}\label{lem:ambient-blowdown}
Let $Y_1\hookrightarrow S$ be an embedding of a smooth projective
surface in a smooth connected projective variety, and let
$b:Y_1\to Y_0$ be the blowdown of a $(-1)$-curve to a smooth projective
surface.  There exist a smooth connected projective variety $S_0$, an
embedding $Y_0\hookrightarrow S_0$, and an isomorphism
$\theta:\pi_1(S)\xrightarrow{\sim}\pi_1(S_0)$ such that
\[
 \theta\circ(\pi_1(Y_1)\to\pi_1(S))
 = (\pi_1(Y_0)\to\pi_1(S_0))\circ b_*.
\]
The equality is understood with compatible basepoints, or up to
conjugation.
\end{lemma}

The proof first places $Y_1$ in a projective threefold and adjusts the
normal bundle of its exceptional curve to
$\mathcal O(-1)\oplus\mathcal O(-1)$.  Blowing up that curve gives an
exceptional $\PP^1\times\PP^1$.  Contracting the other ruling then
performs $b$ on the surface.  This is the classical local two-ruling
construction; compare \cite{Atiyah1958}.  The essential global point
here is to construct a polarization that makes this second contraction
projective.

\begin{proof}
Write $E\subset Y_1$ for the exceptional curve.  Replace $S$ by its
product with a projective space and embed $Y_1$ by the graph of
\[
 Y_1\xrightarrow{b}Y_0\hookrightarrow\PP^r.
\]
The new ambient fundamental group is naturally $\pi_1(S)$.  Let $Q$
be the pullback of $\mathcal O_{\PP^r}(1)$; it is nef, and
$Q|_{Y_1}=b^*L_0$ for a very ample line bundle $L_0$ on $Y_0$.
Further projective-space factors allow us to make the ambient dimension
as large as convenient.

\smallskip
\noindent\emph{A threefold containing the unchanged surface.}
Cut to a smooth projective fourfold $W$ containing $Y_1$, using high
ample degrees.  For each successive cut the rank-loss count
\eqref{proj:rank-loss}, with final dimension four, is at least three;
thus it exceeds $\dim Y_1=2$.  Together with Bertini, this gives smooth
successive cuts, and the Lefschetz theorem preserves the fundamental
group throughout.

For a final threefold, the rank-loss codimension is now two, equal to
the dimension of the surface.  We therefore allow isolated singularities
in the last cut and resolve them while leaving $Y_1$ unchanged.
Take one more sufficiently ample hypersurface $Z\subset W$ through
$Y_1$.  On the surface, its two normal derivatives form a general
section of the rank-two bundle
$N^*_{Y_1/W}\otimes\mathcal O_W(Z)|_{Y_1}$.  In sufficiently high
degree, generation of its first jets makes the zeros transverse and
isolated; since $\dim E=1$, they can also avoid $E$.  Bertini gives
smoothness away from $Y_1$.  At a zero choose coordinates with
$Y_1=(x=y=0)$.  Writing the equation as $xa+yb=0$, transversality says
that $a,b$ restrict to coordinates on $Y_1$.  Hence $x,y,u=a,v=b$ are
coordinates on $W$ and the hypersurface equation is
\begin{equation}\label{proj:node}
 xu+yv=0,\qquad Y_1=(x=y=0).
\end{equation}
Thus $Z$ is a connected normal threefold with only ordinary double
points, all disjoint from $E$.

We require the small resolution of these nodes that leaves $Y_1$
unchanged.  It exists projectively as follows.  Choose a high-degree
Cartier divisor on $Z$ containing $Y_1$, general in its ideal at the
finitely many nodes.  At each node its equation can be taken to be
$x$, after a change of the coordinates in \eqref{proj:node}.  This
divisor is the sum of $Y_1$ and a residual effective Weil divisor $R$.
The latter is Cartier away from the nodes, and at a node its ideal is
$(x,v)$.  Blow up this coherent ideal.  The blowup is projective, is an
isomorphism away from the nodes, and is smooth at every node.  Indeed
its two charts are
\[
 x=tv,\quad y=-tu
 \qquad\text{and}\qquad
 v=sx,\quad u=-sy,
\]
with coordinates $(t,u,v)$ and $(s,x,y)$, respectively.  The strict
transform of $Y_1$ is $t=0$ in the first chart and misses the second
chart.  It therefore maps isomorphically to the original $(u,v)$-plane.
Denote the resulting smooth projective threefold by $T$ and keep the
notation $Y_1\subset T$.

These operations preserve the required fundamental group.  For the
last cut, deform $Z$ to a smooth ample hypersurface in $W$.  Away from
disjoint small neighborhoods of its nodes the family is locally
trivial.  A three-dimensional ordinary double point has simply
connected link and simply connected Milnor fiber, and the small
resolution neighborhood retracts to $\PP^1$.  For completeness, after
writing the node as $\sum_{j=0}^3 z_j^2=0$, its link is the space of
orthonormal two-frames in $\R^4$, hence the unit tangent bundle of
$S^3\simeq\mathrm{SU}(2)$, which is $S^3\times S^2$.  Its smoothing
retracts to $S^3$; writing $z=x+iy$ identifies it with the tangent
bundle of $S^3$ before truncation.  These are also the elementary
quadratic case of the Milnor-fiber theorem \cite{MilnorSingular}.
Van Kampen therefore gives the same fundamental group for the nodal
hypersurface, its smoothing, and its small resolution, compatibly with
the maps to $W$.  Lefschetz for the smoothing yields
$\pi_1(T)\simeq\pi_1(W)\simeq\pi_1(S)$, with the prescribed map
from $\pi_1(Y_1)$.

\smallskip
\noindent\emph{Adjusting the normal bundle.}
The hypersurface degree can be chosen so large that
\[
 N_{Y_1/T}|_E\simeq\mathcal O_E(-m),\qquad m\ge1.
\]
In fact there are no nodes on $E$, and the normal exact sequence there
subtracts $\deg\mathcal O_W(Z)|_E$ from the fixed degree of
$\det N_{Y_1/W}|_E$.  The sequence
\[
 0\longrightarrow\mathcal O_E(-1)
 \longrightarrow N_{E/T}
 \longrightarrow\mathcal O_E(-m)\longrightarrow0
\]
splits because
$H^1(E,\mathcal O_E(m-1))=0$.  Blow up the threefold along $E$.
Since $E$ is a Cartier divisor on $Y_1$, the strict transform of $Y_1$
is isomorphic to $Y_1$.  Its normal line bundle becomes
$N_{Y_1/T}\otimes\mathcal O_{Y_1}(-E)$, so its degree on $E$ increases
by one, because $E^2=-1$.  Repeating $m-1$ times gives a smooth
projective threefold, still denoted $T$, in which
\[
 N_{E/T}\simeq\mathcal O_E(-1)\oplus\mathcal O_E(-1),
\]
with the first summand the tangent-normal direction inside $Y_1$.
At each step the same extension vanishing justifies this splitting.

Blow up $E$ once more, obtaining $\widehat T$, and let $\widehat Y$
be the strict transform of $Y_1$.  The exceptional divisor is
\begin{equation}\label{proj:exceptional-quadric}
 G=\PP^1_{\mathrm{curve}}\times\PP^1_{\mathrm{normal}},\qquad
 N_{G/\widehat T}=\mathcal O_G(-1,-1),\qquad
 \widehat Y\cap G=\PP^1_{\mathrm{curve}}\times\{q\}.
\end{equation}
The intersection is transverse.  Moreover
$\mathcal O_{\widehat T}(\widehat Y)|_G=\mathcal O_G(0,1)$, and
$N_{\widehat Y/\widehat T}|_E=\mathcal O_E$.
We continue to denote by $Q$ the pullback of the earlier nef bundle;
it is trivial on $G$ and restricts on $\widehat Y\simeq Y_1$ to
$b^*L_0$.  Figure~\ref{fig:blowdown} shows the two rulings and the
curve on the surface that the second contraction must remove.

\begin{figure}[htbp]
\centering
\begingroup
\definecolor{blowaccent}{RGB}{22,91,119}
\begin{tikzpicture}[x=1cm,y=1cm,>=Latex,
  every node/.style={font=\small},
  marked/.style={blowaccent,line width=1.5pt}]
  \coordinate (sw) at (3.6,.65);
  \coordinate (se) at (7.8,.65);
  \coordinate (nw) at (4.5,3.25);
  \coordinate (ne) at (8.7,3.25);
  \fill[black!3] (sw)--(se)--(ne)--(nw)--cycle;
  \foreach \t in {.25,.5,.75}{
    \draw[black!23] ($(sw)!\t!(nw)$)--($(se)!\t!(ne)$);
    \draw[black!23] ($(sw)!\t!(se)$)--($(nw)!\t!(ne)$);
  }
  \draw[black!55] (sw)--(se)--(ne)--(nw)--cycle;
  \coordinate (el) at ($(sw)!.5!(nw)$);
  \coordinate (er) at ($(se)!.5!(ne)$);
  \draw[marked] (el)--(er);
  \node[above=6pt] at (6.65,3.25)
    {$G=\mathbb P^1_{\mathrm{curve}}\times\mathbb P^1_{\mathrm{normal}}$};
  \node[fill=white,inner sep=2pt,text=blowaccent] at (6.15,1.95)
    {$\widehat Y\cap G=E\times\{q\}$};
  \node[below=4pt] at (5.7,.65) {curve direction};

  \draw[marked] (0,.7)--(2.05,.7);
  \node[below=5pt] at (1.025,.7) {$E\subset Y_1\subset T$};
  \draw[->,line width=.75pt] (3.45,1.65)--(1.65,1.05);
  \node[align=center] at (1.6,2.35)
    {original blowdown\\$G\longrightarrow\mathbb P^1_{\mathrm{curve}}$};

  \draw[black!65,line width=.8pt] (11.1,.75)--(11.7,2.65);
  \fill[blowaccent] (11.4,1.7) circle (2.8pt);
  \node[right=5pt,align=left] at (11.4,1.7) {$q$\\in $Y_0$};
  \draw[->,line width=.75pt] (8.7,1.9)--(10.7,1.9);
  \node[align=center] at (10.35,3.65)
    {other blowdown\\$G\longrightarrow\mathbb P^1_{\mathrm{normal}}$};
  \node[below=6pt,align=center] at (11.1,.6)
    {image curve\\in $T'$};
\end{tikzpicture}
\endgroup
\caption{The two rulings of the last exceptional divisor in
  $\widehat T$. The original blowdown contracts the normal-direction
  fibers and leaves the curve $E$ in $Y_1$ unchanged. The other
  blowdown, whose target threefold $T'$ is constructed below,
  contracts the horizontal ruling, including
  $\widehat Y\cap G=E\times\{q\}$, so the surface undergoes its
  point blowdown $Y_1\to Y_0$. The parallelogram records only the
  product structure of $G$, whose normal bundle is
  $\mathcal O_G(-1,-1)$.}
\label{fig:blowdown}
\end{figure}


\smallskip
\noindent\emph{A polarization for the other contraction.}
We seek an ample line bundle $H$ restricting to
$\mathcal O_G(\alpha,\beta)$, with $\beta>\alpha>0$.
Since $\mathcal O_{\widehat T}(G)|_G=\mathcal O_G(-1,-1)$,
adding $\alpha G$ will then give the restriction
$\mathcal O_G(0,\beta-\alpha)$: it is trivial on the curves to be
contracted and positive on the retained factor.
Choose an ample integral line bundle $H_0$ on $\widehat T$, with
$H_0|_G=\mathcal O_G(\alpha,\beta_0)$, where $\alpha,\beta_0>0$.
Adding $k\widehat Y$ increases the second degree by $k$ and leaves the
first unchanged.  Choose $k\ge0$ with $\beta:=\beta_0+k>\alpha$.
On $\widehat Y$, the line bundle $H_0+k\widehat Y$ still has positive
degree on $E$, because $N_{\widehat Y/\widehat T}|_E$ is trivial.
It is therefore relatively ample for $b:\widehat Y\to Y_0$, whose
only positive-dimensional fiber is $E$.  Adding a sufficiently large
multiple $lQ$ makes its restriction to $\widehat Y$ ample.

We claim that
\[
 H=H_0+k\widehat Y+lQ
\]
is ample on all of $\widehat T$.  Here $H_0+lQ$ is ample since $Q$
is nef.  We use the following consequence of the Nakai--Moishezon
criterion: if $A_0$ is ample, $D_0$ is an effective Cartier divisor,
$k\ge0$, and $(A_0+kD_0)|_{D_0}$ is ample, then $A_0+kD_0$ is ample.
Indeed, for an irreducible $r$-dimensional subvariety $V$ not contained
in $D_0$, the difference of top intersections is
\[
 \bigl((A_0+kD_0)^r-A_0^r\bigr)\cdot V
 = kD_0\cdot V\cdot
   \sum_{j=0}^{r-1}(A_0+kD_0)^j A_0^{r-1-j}\ge0.
\]
Every term on the right is a mixed intersection of ample restrictions
on the effective cycle $D_0\cap V$.  For $V\subset D_0$, positivity
follows directly from the restriction hypothesis.  Nakai--Moishezon
now applies.  Taking $D_0=\widehat Y$ proves the claim.  We have obtained
\begin{equation}\label{proj:ample-degrees}
 H|_G=\mathcal O_G(\alpha,\beta),\qquad \beta>\alpha>0.
\end{equation}
The relative ampleness, Serre vanishing, and numerical ampleness facts
used here are standard projective criteria; see \cite{Hartshorne}.

Set $D=H+\alpha G$.  We shall prove directly that a high multiple of
$D$ contracts $G$ to its second factor and has smooth projective image.
For a sufficiently large integer $n$, put
$L_j=nD-jG$.  Then
\[
 L_{n\alpha}=nH,
 \qquad L_j|_G=\mathcal O_G\bigl(j,n(\beta-\alpha)+j\bigr).
\]
Serre vanishing gives $H^1(\widehat T,L_{n\alpha})=0$.
For $0\le j<n\alpha$, the restriction exact sequences and
$H^1(G,\mathcal O_G(j,n(\beta-\alpha)+j))=0$ propagate this
vanishing downwards:
\begin{equation}\label{proj:vanishing}
 H^1(\widehat T,L_j)=0\qquad(0\le j\le n\alpha).
\end{equation}
Choose also $n\alpha\ge2$ and $nH$ very ample.  In particular the
restriction maps for $L_0$ and $L_1$ are surjective.  The former gives
all sections of $\mathcal O_G(0,n(\beta-\alpha))$.  Hence $|nD|$ is
basepoint free on $G$ and restricts to the second projection followed
by a Veronese embedding.  Away from $G$, it contains the system
obtained by multiplying sections of $nH$ by the canonical section of
$n\alpha G$.  This subsystem embeds the complement of $G$ and
separates each of its points from $G$.  Thus $|nD|$ defines a projective
morphism
\[
 \varphi:\widehat T\longrightarrow T'
\]
that is an isomorphism off $G$ and whose nontrivial fibers are exactly
the curves $\PP^1_{\mathrm{curve}}\times\{t\}$.

\smallskip
\noindent\emph{Smoothness of the contracted space and of the surface.}
We check the local model, rather than assume that an analytic
contraction is projective or smooth.  Fix a fiber
$F_t=\PP^1_{\mathrm{curve}}\times\{t\}$ of $G\to\PP^1_{\mathrm{normal}}$.
Use a section of $nD$ nonzero on $F_t$ as an affine denominator.
The restriction surjection for $L_0$ gives a section ratio $z$ whose
restriction to $G$ is a local coordinate on its second factor.
The restriction surjection for $L_1=nD-G$ gives two section ratios
$x,y$ vanishing on $G$, whose first normal coefficients on $F_t$ are
a basis of $H^0(\PP^1,\mathcal O(1))$.

If $e=0$ is a local equation for $G$, write $x=ea$, $y=eb$.
The functions $a,b$ have no common zero near $F_t$.  Thus
$[x:y]=[a:b]$ extends across $G$, and the map $(z,x,y)$ lifts to the
blowup of $\C^3$ along the axis $x=y=0$.  On $F_t$, its projective
coordinate is the standard isomorphism to the exceptional $\PP^1$.
In the chart $a\ne0$, the coordinates $(z,x,y/x)$ have nonsingular
Jacobian along the fiber: they respectively detect its base direction,
the direction normal to $G$, and the tangent direction of the fiber.
The same holds in the other chart.  If injectivity failed on every
smaller neighborhood of $F_t$, colliding pairs would have subsequences
converging to two points of this compact fiber.  Injectivity on the fiber
identifies their limits, contradicting local biholomorphism there; hence
the lifted map is an isomorphism of neighborhoods.  Properness of the
standard blowup then allows us to shrink its image to the full inverse
image of a neighborhood of the point on the axis.

All the other affine section ratios defining $\varphi$ are holomorphic
on this blowup neighborhood and descend to holomorphic functions on
the unblown-up neighborhood.  One may use here the elementary identity
$\pi_*\mathcal O_{\operatorname{Bl}_{\mathrm{axis}}\C^3}
=\mathcal O_{\C^3}$: a function descends off the codimension-two axis
and extends across it by normality.  Since the entire fiber of
$\varphi$ over $\varphi(F_t)$ is $F_t$, properness ensures that the
resulting neighborhood computes the germ of its projective image.
In these coordinates that image is a graph over $(z,x,y)$.
Consequently $T'$ is smooth, and $\varphi$ is the ordinary smooth
blowdown of $G$ along its other ruling.

At $F_q=\widehat Y\cap G$, transversality in
\eqref{proj:exceptional-quadric} makes projection to the blowup of the
$(x,y)$-plane a local isomorphism on $\widehat Y$ along $F_q$, and it
is an isomorphism on that fiber.  The same compactness and shrinking
argument gives an isomorphism of neighborhoods with the surface
blowup.  Its remaining coordinate $z$ is a holomorphic function there and
descends to the plane, by the same extension argument.  Thus the
image of $\widehat Y$ is a smooth surface, and its restriction is
exactly the point blowdown of $E$.  By uniqueness of that blowdown,
this image identifies with $Y_0$.

Finally, smooth blowups and blowdowns along smooth centers of complex
codimension at least two preserve fundamental groups.  This follows
from the usual tubular-neighborhood model and van Kampen, or from
general position together with the simply connected projective-space
fibers.  The earlier threefold construction preserved the groups
compatibly with the surface inclusion.  The final diagram
\[
 \widehat Y\longrightarrow\widehat T\longrightarrow T',
 \qquad \widehat Y\simeq Y_1\xrightarrow{b}Y_0\subset T'
\]
shows that the induced identifications also preserve the homomorphism
from the surface.  Taking $S_0=T'$ gives the required $\theta$.
\end{proof}

\subsection{Completion of the construction}

\begin{theorem}\label{thm:cyclic-embedding}
There exist a simple abelian surface $A$, a smooth connected projective
complex variety $S'$, and a closed embedding $A\hookrightarrow S'$
such that
\[
 \im\bigl(\pi_1(A)\longrightarrow\pi_1(S')\bigr)\simeq\Z.
\]
\end{theorem}

\begin{proof}
Theorem~\ref{thm:orbifold-cyclic} gives an iterated point blowup $Y$ of
the simple abelian surface $A$, a smooth connected effective proper
orbifold $\cS$ with projective coarse space, and a map $Y\to\cS$
whose image lies in the ordinary locus and whose fundamental-group
image is infinite cyclic.  Lemma~\ref{lem:ambient-realization} replaces
this map by an embedding $Y\hookrightarrow S$ with the same group
image.  Apply Lemma~\ref{lem:ambient-blowdown} to the point blowups in
reverse order.  Every step preserves both the ambient fundamental
group and the homomorphism from the surface.  After the last step the
surface is $A$, giving the asserted embedding and cyclic image.
\end{proof}

\begin{samepage}
\begin{proof}[Proof of Theorem~\ref{thm:main}]
Apply Proposition~\ref{prop:large-reduction} to the embedding supplied
by Theorem~\ref{thm:cyclic-embedding}.  The resulting smooth connected
projective fourfold contains $A$ with infinite cyclic fundamental-group
image, has large fundamental group, and has a non-Stein universal cover.  By Lemma~\ref{red:compact-criterion}, its universal
cover contains no positive-dimensional compact complex-analytic
subvariety.
\end{proof}
\end{samepage}
